\documentclass[11pt,a4paper]{amsart}
\usepackage{amsmath,amssymb,amsthm,mathtools}
\usepackage[margin=1.2in]{geometry}
\usepackage{hyperref}

\theoremstyle{plain}
\newtheorem{theorem}{Theorem}[section]
\newtheorem{lemma}[theorem]{Lemma}
\newtheorem{proposition}[theorem]{Proposition}
\newtheorem{corollary}[theorem]{Corollary}
\theoremstyle{definition}
\newtheorem{definition}[theorem]{Definition}
\newtheorem{conjecture}[theorem]{Conjecture}
\theoremstyle{remark}
\newtheorem{observation}[theorem]{Observation}
\newtheorem{remark}[theorem]{Remark}

\newcommand{\ZZ}{\mathbb{Z}}
\newcommand{\QQ}{\mathbb{Q}}
\newcommand{\FF}{\mathbb{F}}
\newcommand{\CC}{\mathbb{C}}
\newcommand{\vp}{v_p}
\newcommand{\A}{A}
\DeclareMathOperator{\Res}{Res}

\title{A two-level digit law for Ap\'ery's pair}
\author{[Author]}
\date{}

\begin{document}
\begin{abstract}
Let $a_n$ be the Ap\'ery numbers for $\zeta(3)$ and $b_n$ their companion. We show that
the pair $(a_n,\,p^3b_n)$ transforms under base-$p$ digit splitting by a single
upper-triangular matrix whose diagonal is one explicit scalar and whose off-diagonal
entry is the low digit's own companion value, scaled by exactly $p^{3}$. The scalar is
the second-order truncation of a $\Gamma$-deformation of the summand, summed over the
full digit range. We prove the diagonal of the law to second order, for both rows, and
with it the rank-one profile of the digit defect. Everything reduces to the residue
theorem for one explicit rational function, applied twice: once over $\CC$, giving an
identity over $\QQ$, and once over $\FF_p$, where the carry region is exactly the
complement of the pole divisor. Only the order-$p^4$ off-diagonal remains conjectural.
\end{abstract}
\maketitle

Throughout $p$ is a prime $\ge5$, and for $x,y\in\QQ$ we write $x\equiv y\pmod{p^m}$ to
mean $\vp(x-y)\ge m$. Evidence is labelled: statements called Theorem, Lemma,
Proposition or Corollary are proved; Conjectures are not, and the single Conjecture
below records exactly what is not.

\section{The pair, and the two functionals}

\begin{definition}
$\A(n,k)=\binom nk^2\binom{n+k}k^2$, \ $a_n=\sum_{k=0}^n \A(n,k)$, \ $H^{(r)}_m=\sum_{j=1}^m j^{-r}$,
$H_m=H^{(1)}_m$, and $b_n=\sum_{k=0}^n \A(n,k)\bigl(2H^{(3)}_n-H^{(3)}_k\bigr)$, the solution of
\[
  (n+2)^3b_{n+2}=(34n^3+153n^2+231n+117)\,b_{n+1}-(n+1)^3b_n,\qquad b_0=0,\ b_1=6 .
\]
\end{definition}

\begin{definition}\label{def:uv}
$u(r,s)=H_{r+s}-H_{r-s}$ and $v(r,s)=H_{r+s}+H_{r-s}-2H_s$; these are
$\tfrac12\partial_n\log \A$ and $\tfrac12\partial_k\log \A$. Put
\[
  U_r=\sum_{s} \A(r,s)\,u(r,s),\qquad V_r=\sum_{s} \A(r,s)\,v(r,s).
\]
\end{definition}

$U_0,\dots,U_8=0,\;6,\;105,\;2219,\;\tfrac{104825}2,\;\tfrac{13276637}{10},\;
\tfrac{70543291}2,\;\tfrac{67890874657}{70},\;\tfrac{766399019471}{28}$.
Note $U_r=\tfrac12\,\mathrm{d}a_n/\mathrm{d}n$ in the $\Gamma$-form.

\section{The expansion, and why the defect is rank one}

\begin{lemma}[digit expansion]\label{lem:exp}
For $0\le s\le r<p$, $0\le c\le a<p$ and $r+s<p$,
\[
  \A(ap+r,\,cp+s)\;\equiv\;\A(a,c)\,\A(r,s)\,\bigl(1+2p[\,a\,u(r,s)+c\,v(r,s)\,]\bigr)\pmod{p^2}.
\]
Outside that region the relevant binomial is $\equiv0\pmod p$; since $\A$ is a square,
both $\A(ap+r,cp+s)$ and $\A(r,s)$ are then $\equiv0\pmod{p^2}$, so those regions do not
contribute at first order at all.
\end{lemma}

\begin{proof}
Only the weak Wolstenholme congruence $H_{p-1}\equiv0\pmod p$ is needed (the $(ip)^{2}$
terms die of their own accord), and it follows from pairing $l\leftrightarrow p-l$. It
makes the $(ap)!/(p^aa!)\equiv((p-1)!)^a$ factors cancel, giving
$\binom{ap+r}{cp+s}\equiv\binom ac\binom rs(1+p[aH_r-cH_s-(a-c)H_{r-s}])$ and
$\binom{(a+c)p+(r+s)}{cp+s}\equiv\binom{a+c}c\binom{r+s}s(1+p[(a+c)H_{r+s}-cH_s-aH_r])$
modulo $p^2$. In the product $\A$ the two $aH_r$ terms cancel.
\end{proof}

Lemma \ref{lem:exp} is available only on $r+s<p$, so summing it produces the
\emph{restricted} functionals $\widetilde U_r=\sum_{s\le r,\ r+s<p}\A(r,s)u(r,s)$ and
$\widetilde V_r$, not $U_r$ and $V_r$. The two agree:

\begin{lemma}\label{lem:restrict}
For $r<p$, \ $\widetilde U_r\equiv U_r$ and $\widetilde V_r\equiv V_r \pmod p$.
\end{lemma}

\begin{proof}
For $r+s\ge p$ the carried binomial satisfies $\binom{r+s}s\equiv0\pmod p$, and $\A$ is a
square, so $p^2\mid \A(r,s)$; while $u(r,s),v(r,s)$ have arguments $<2p$ and hence
$\vp\ge-1$. Every dropped term therefore has $\vp\ge1$.
\end{proof}

This is not a formality: $U_r$ is a $p$-independent rational which \emph{does} carry $p$ in
its denominator for small $r$ — $U_5=13276637/10$ and $U_7=67890874657/70$, so
$v_5(U_5)=v_7(U_7)=-1$. It is $p$-integral precisely because $r<p$, which is a consequence
of Lemma \ref{lem:restrict} rather than of anything stated before it.

Since $H^{(3)}_{ap+r}=p^{-3}H^{(3)}_a+G_{ap+r}$ with $G_m=\sum_{j\le m,\ p\nmid j}j^{-3}$
$p$-integral, the weight contributes nothing below order $p^3$. Lemma \ref{lem:exp}
therefore gives, with $a'_a=\sum_c c\,\A(a,c)$ and $b'_a=\sum_c c\,\A(a,c)(2H^{(3)}_a-H^{(3)}_c)$,
\[
  \tfrac1p\bigl(a_{ap+r}-a_aa_r\bigr)\equiv 2\bigl[a\,a_aU_r+a'_aV_r\bigr],
  \qquad
  \tfrac1p\bigl(p^3b_{ap+r}-b_aa_r\bigr)\equiv 2\bigl[a\,b_aU_r+b'_aV_r\bigr]
  \pmod p,
\]
a form of rank at most two. One channel degenerates identically:

\begin{theorem}\label{thm:V}
$V_n=\displaystyle\sum_{k=0}^{n}\binom nk^2\binom{n+k}k^2\bigl(H_{n+k}+H_{n-k}-2H_k\bigr)=0$
for all $n\ge0$.
\end{theorem}

\begin{proof}
$V_n=\tfrac12\sum_k\partial_k \A_\Gamma(n,k)$ with
$\A_\Gamma(n,z)=[\Gamma(n+z+1)/(\Gamma(z+1)^2\Gamma(n-z+1))]^2$, and reflection gives
$\A_\Gamma(n,z)=(\sin^2\pi z/\pi^2)\,g(z)$ with
$g(z)=\Gamma(n+z+1)^2\Gamma(z-n)^2/\Gamma(z+1)^4$. No analysis is required, because $g$ is
\emph{rational}:
\[
  g=\varphi^2,\qquad
  \varphi(z)=\frac{\prod_{i=1}^{n}(z+i)}{\prod_{j=0}^{n}(z-j)} .
\]
At an integer $z=k$ the double poles of $\A_\Gamma$ cancel and
$\partial_z \A_\Gamma|_{z=k}=\Res_{z=k}g$. Now $\varphi$ has simple poles with residues
$\beta_k=(-1)^k\binom nk\binom{n+k}k$ — this is Lagrange interpolation of
$\prod_{i=1}^n(X+i)$ at the nodes $0,\dots,n$ — so
\[
  \Res_{z=m}\varphi^2=2\beta_m\sum_{k\neq m}\frac{\beta_k}{m-k},
\]
which is antisymmetric in $(m,k)$ and therefore sums to zero over $m$.
\end{proof}

\begin{remark}
No telescoping proof of Theorem \ref{thm:V} of the shape used for the closed form exists:
any $\Psi=\A\cdot(\alpha H_{n+k}+\beta H_{n-k}+\gamma H_k+\delta)$ forces
$r(k)\alpha(k+1)-\alpha(k)=1$, which is Gosper's equation for the Ap\'ery summand.
\end{remark}

\begin{corollary}\label{cor:first}
For $a,r<p$,
\[
  \tfrac1p\bigl(a_{ap+r}-a_aa_r\bigr)\equiv 2a\,a_a\,U_r,
  \qquad
  \tfrac1p\bigl(p^3b_{ap+r}-b_aa_r\bigr)\equiv 2a\,b_a\,U_r \pmod p .
\]
Both rows carry the same functional $U_r$ and differ only by their own value at level $a$.
Equivalently
\[
  (a_n,\;p^3b_n)\;\equiv\;(a_a,\;b_a)\cdot\bigl(a_r+2p\,a\,U_r\bigr)\pmod{p^2},
  \qquad n=ap+r .
\]
\end{corollary}

\begin{remark}
$U_r$ is not $a_r$, nor $b_r$, nor divisible by $a_r$: at $p=5$, $U_r\not\equiv0$ exactly
where $a_r\equiv0$. It is the exact analogue, for this pair, of the functional appearing
in the corresponding law for the Brown--Zudilin $\zeta(5)$ family, and has the same
origin, $\partial_n\log$ of the summand.
\end{remark}

\section{The scalar, and the matrix}

\begin{definition}\label{def:adig}
With $\Pi(t)=\prod_{i\le t}(1+\varepsilon/i)$ and $\bar\Pi(t)=\prod_{i\le t}(1-\varepsilon/i)$,
\[
  \A_\Gamma(r+\varepsilon,s)=\A(r,s)\frac{\Pi(r+s)^2}{\Pi(r-s)^2}\quad(s\le r),
\]
\[
  \A_\Gamma(r+\varepsilon,r+m)=\varepsilon^2 C(r,m)\,\Pi(2r+m)^2\bar\Pi(m-1)^2
  \quad(m\ge1),
\]
where $C(r,m):=\bigl((2r+m)!\,(m-1)!/((r+m)!)^2\bigr)^{2}$,
and $\operatorname{Adig}(p,r;\varepsilon)=\sum_{s=0}^{p-1}\A_\Gamma(r+\varepsilon,s)$, the sum
being over the \emph{full} digit range. Write $X_p(r)=[\varepsilon^2]\operatorname{Adig}(p,r)$.
\end{definition}

Then $[\varepsilon^0]=a_r$ and $[\varepsilon^1]=2U_r$. The terms with digit $s>r$ — the
borrow region — are present and are exactly of order $\varepsilon^2$.

\begin{theorem}[the two-level digit law]\label{thm:scalar}
Let $p\ge5$ be prime and $n=ap+r$ with $0\le a,r<p$. With
$u(a,r)=a_r+2p\,a\,U_r+p^2a^2X_p(r)$,
\[
  (a_n,\;p^3b_n)\;\equiv\;(a_a,\;b_a)\cdot u(a,r) \pmod{p^3};
\]
equivalently $u(a,r)=\bigl[\sum_{s=0}^{p-1}\A_\Gamma(r+\varepsilon,s)\bigr]_{\varepsilon=pa}$
truncated at order $\varepsilon^2$. The floor is exactly $3$.
\end{theorem}

The proof is \S\ref{sec:second}. Truncating $\operatorname{Adig}$ at orders $m=0,1,2$
gives floors $1,2,3$, and floor $3$ for every $m\ge3$: the scalar picture saturates at
depth three. The restricted deformation $\sum_{s\le r}$ saturates at depth two, so the
borrow region is exactly the second-order correction.

\begin{corollary}\label{cor:rank}
The digit defect has rank $\le1$ at first order and rank $\le1$ at second order, for both
rows, with $a$-side factor exactly $a\cdot(a_a\mid b_a)$ respectively
$a^2\cdot(a_a\mid b_a)$, and with the \emph{same} $r$-side vector for both rows:
$(U_r)_r$ at first order, $(X_p(r))_r$ at second. The rank is exactly $1$ iff the
respective $r$-vector is $\not\equiv0$, which holds at all 13 primes $5\le p\le47$.
\end{corollary}

\begin{proposition}[the cross term]\label{prop:cross}
Modulo $p$ a full block satisfies $\sum_{t=1}^{p-1}t^{-3}\equiv0$, so
$G_{ap+r}\equiv H^{(3)}_r$ and $G_{cp+s}\equiv H^{(3)}_s$. Hence the piece of $p^3b_n$
neglected above is
\[
  p^3\sum_k \A(n,k)\bigl(2G_n-G_k\bigr)\;\equiv\;p^3\,a_a\,b_r \pmod {p^4}.
\]
\end{proposition}

That is, the low digit acts on the pair by

\begin{conjecture}\label{conj:matrix}
For every prime $p\ge5$ and $n=ap+r$ with $a,r<p$,
\[
  (a_n,\;p^3b_n)\;\equiv\;(a_a,\;b_a)\cdot
  \begin{pmatrix} u(a,r) & p^3\,b_r\\[2pt] 0 & u(a,r)\end{pmatrix}
  \pmod{p^4},
\]
with $u(a,r)$ as in Theorem \ref{thm:scalar}.
\end{conjecture}

Its diagonal is now proved: the $\bmod\ p^2$ truncation is Corollary \ref{cor:first} and
the $\bmod\ p^3$ truncation is Theorem \ref{thm:scalar}. What remains conjectural is only
the off-diagonal entry at order $p^4$. That entry is derived in Proposition
\ref{prop:cross} and verified at 13 primes in the sharp form: at order $p^3$ the $b$-row
defect has rank $3$, subtracting $p^3a_ab_r$ drops it to rank $2$, and its $r$-side row
space then coincides with the $a$-row's. The residual scalar defect at $p^3$ has rank $2$
and is not yet named; that alone is why the statement is still a conjecture.

\begin{remark}
The cross entry is the low digit's own companion value, scaled by exactly $p^{3}=p^{w}$
where $w=3$ is the weight of $\zeta(3)$. So the weight is precisely the order at which
the off-diagonal switches on. For the rank-three Brown--Zudilin $\zeta(5)$ family the
corresponding statement is that a cross term cannot appear before order $p^3$; here it
is derived rather than bounded, and the entry is named.
\end{remark}

\section{The second order}\label{sec:second}

Writing $\A(ap+r,cp+s)\equiv \A(a,c)\A(r,s)\exp(2p\Lambda_1-p^2\Lambda_2)$ with
$\Lambda_1=a\,u+c\,v$ and
$\Lambda_2=(a+c)^2H^{(2)}_{r+s}-2c^2H^{(2)}_s-(a-c)^2H^{(2)}_{r-s}$, the bracket
$2\Lambda_1^2-\Lambda_2$ has three channels, and we write $\Sigma_{a^2},\Sigma_{ac},
\Sigma_{c^2}$ for their $\A(r,s)$-weighted sums over $s\le r$:
\[
  a^2:\ 2u^2-\bigl(H^{(2)}_{r+s}-H^{(2)}_{r-s}\bigr),\qquad
  ac:\ 4uv-2\bigl(H^{(2)}_{r+s}+H^{(2)}_{r-s}\bigr),
\]
\[
  c^2:\ 2v^2-\bigl(H^{(2)}_{r+s}-2H^{(2)}_s-H^{(2)}_{r-s}\bigr).
\]
Theorem \ref{thm:scalar} is equivalent to the vanishing of the last two modulo $p$. They
are \emph{not} rational identities: $\Sigma_{ac}=0,\,-26,\,-\tfrac{905}2,\,
-\tfrac{167965}{18},\dots$ and $\Sigma_{c^2}=0,\,13,\,\tfrac{905}4,\,
\tfrac{167965}{36},\dots$, nonzero over $\QQ$ for $1\le r\le30$. Both statements follow
from one rational function.

\begin{lemma}[reflection, rational form]\label{lem:grat}
$\A_\Gamma(r,z)=\dfrac{\sin^2\pi z}{\pi^2}\,g_r(z)$ with
$g_r(z)=\Bigl[\prod_{i=1}^{r}(z+i)\big/\prod_{j=0}^{r}(z-j)\Bigr]^2$, a rational
function of degree $-2$ whose only poles are double poles at $z=0,1,\dots,r$.
\end{lemma}

This replaces the $\Gamma$-asymptotics of Theorem \ref{thm:V} by an identity. Since
$\sin^2\pi z/\pi^2=(z-s)^2-\tfrac{\pi^2}3(z-s)^4+O((z-s)^6)$ at an integer $s$, a
function $f=\alpha/(z-s)^2+\beta/(z-s)+\gamma+O(z-s)$ gives
$\Phi=(\sin^2\pi z/\pi^2)f$ regular at $s$ with $\Phi(s)=\alpha$, $\Phi'(s)=\beta$ and
$\tfrac12\Phi''(s)=\gamma-\tfrac{\pi^2}3\alpha$. For $f=g_r$ this reads
\[
  \alpha_s=\A(r,s),\qquad \beta_s=2\A(r,s)\,v(r,s),\qquad
  \gamma_s=\A(r,s)\bigl(2v^2-(H^{(2)}_{r+s}-2H^{(2)}_s-H^{(2)}_{r-s})\bigr),
\]
so $\sum_s\beta_s=2V_r$ and $\sum_s\gamma_s=\Sigma_{c^2}(r)$; and for $m\ge1$,
\[
  g_r(r+m)=C(r,m),
\]
which is exactly the $\varepsilon^2$ weight of the borrow digit $s=r+m$. The $\pi^2/3$
of the dictionary cancels the $\zeta(2)$ of $\partial_k^2\log\A_\Gamma$, which is why
$\gamma_s$ is the clean $c^2$ channel.

\begin{theorem}\label{thm:R1}
$\Sigma_{ac}(r)+2\,\Sigma_{c^2}(r)=0$ for every $r\ge0$, as an identity over $\QQ$.
\end{theorem}

\begin{proof}
$\Sigma_{ac}+2\Sigma_{c^2}=\sum_{s\le r}\partial_k\bigl[(\partial_n+\partial_k)
\A_\Gamma\bigr](r,s)$; the $\zeta(2)$ terms enter with weights $4$ and $2\cdot(-2)$ and
cancel, so this is the unique $\zeta(2)$-free combination. Now $\sin^2\pi(z-n)$ is
invariant along $(n,z)\mapsto(n+t,z+t)$, and in that direction
$(\partial_n+\partial_z)\log g_n=4[\psi(n+z+1)-\psi(z+1)]$ — the $\psi(z-n)$ cancels.
Hence at $n=r\in\ZZ_{\ge0}$
\[
  h_r(z):=(\partial_n+\partial_z)g_n(z)\big|_{n=r}=4\,g_r(z)\sum_{i=1}^{r}\frac1{z+i}
\]
is rational, with poles only the double poles at $z=0,\dots,r$ (at $z=-i$ the simple
pole meets the double zero of $g_r$), and $h_r=O(z^{-3})$. A rational function that is
$O(z^{-2})$ at infinity has vanishing residue sum, and by the dictionary
$\partial_z[(\partial_n+\partial_z)\A_\Gamma](r,z)|_{z=s}=\operatorname*{Res}_{z=s}h_r$.
\end{proof}

\begin{lemma}\label{lem:fp}
Let $p\ge5$, $0\le r<p$, and let $f\in\QQ(z)$ have poles only at $0,\dots,r$, of order
$\le2$, with $p$-integral principal parts and $\deg f\le-2$. Then
$\sum_{s=0}^{r}\operatorname{FP}_{z=s}f+\sum_{t=r+1}^{p-1}f(t)\equiv0\pmod p$.
\end{lemma}

\begin{proof}
Both sums are sums of $\alpha_{s'}/(y-s')^2+\beta_{s'}/(y-s')$ over $y\ne s'$ with
$0\le y\le p-1$, all terms $p$-integral. Interchanging, for fixed $s'$ the difference
$y-s'$ runs over all nonzero residues, and $\sum_{w\in\FF_p^\times}w^{-1}
=\sum_{w\in\FF_p^\times}w^{-2}\equiv0$; the second needs $p\ge5$.
\end{proof}

\begin{theorem}\label{thm:R2}
For $p\ge5$ and $0\le r<p$, $\ \Sigma_{c^2}(r)+\Xi_p(r)\equiv0\pmod p$, where
$\Xi_p(r)=\sum_{m=1}^{p-1-r}C(r,m)$.
\end{theorem}

\begin{proof}
Take $f=g_r$ in Lemma \ref{lem:fp}: $\deg g_r=-2$, the principal parts are $p$-integral
($\alpha_s=\A(r,s)\in\ZZ$, and the only $p$ in a denominator of $\beta_s$ is the $1/p$ of
$H_{r+s}$ when $r+s\ge p$, where $v_p(\A(r,s))=2$), $\sum_s\gamma_s=\Sigma_{c^2}(r)$, and
$\sum_{t=r+1}^{p-1}g_r(t)=\Xi_p(r)$.
\end{proof}

Theorem \ref{thm:scalar} follows, with $X_p(r)=\Sigma_{a^2}(r)+\Xi_p(r)$.

\begin{remark}
The pole set of $g_r$ is $\{0,\dots,r\}$ — the region-I digit range — and its complement
in $\FF_p$ is $\{r+1,\dots,p-1\}$, the borrow range. So Theorem \ref{thm:R2} says the
finite parts on the pole divisor cancel the values off it, over $\FF_p$: it is the
residue theorem read over $\FF_p$ rather than over $\CC$. That is why the vanishing
cannot be a rational identity, and indeed cannot be: $\Sigma_{c^2}(r)$ is a fixed nonzero
rational for $r\ge1$ while $\Xi_p(r)$ depends on $p$. Dropping the borrow region
($\Xi=0$) breaks Theorem \ref{thm:R2} at every $p$ and every $r$ with
$\Sigma_{c^2}(r)\not\equiv0$. The contrast with the first order is that there the borrow
region contributes nothing at all, because $\A$ is a square (Lemma \ref{lem:exp}).
\end{remark}

\begin{remark}[$p=2,3$]
Both hypotheses $p\ge5$ are used: Lemma \ref{lem:fp} needs
$\sum_{\FF_p^\times}w^{-2}\equiv0$, and the digit expansion needs
$W(M)\equiv((p-1)!)^M \bmod p^3$, i.e.\ Wolstenholme. At $p=2,3$ both fail, as do the
regional formulas and Theorem \ref{thm:R2} at every $r$; only Theorem \ref{thm:R1},
being $p$-free, survives. The conclusion of Theorem \ref{thm:scalar} nevertheless holds
at $p=2$ and $p=3$ over all $4$ and $12$ admissible cells. This is a resolution artefact
rather than a theorem: the statement is equivalent to
$a\,m_1(\Sigma_{ac}-2\Xi)+m_2(\Sigma_{c^2}+\Xi)\equiv0$, whose necessity requires the
vectors $(a\,m_1(a),m_2(a))_{1\le a<p}$ to span $\FF_p^2$, and their rank is $0$ at
$p=2$ and $1$ at $p=3$ against $2$ at every $p\ge5$. With one usable digit, and two, the
obstruction cannot be seen.
\end{remark}

\end{document}
